\section{第 4 章：算术表达式的 ML 实现}

\begin{taplauthorsolutionentrybox}
\phantomsection
\label{sol:author:4.2.1}
\textbf{4.2.1 解答}

\texttt{eval} 每次递归调用自身时，都会在调用栈上再压入一个
\texttt{try} 处理器。每次单步求值都会产生一次递归调用，所以调用栈最终会
溢出。实质上，用 \texttt{try} 包住对 \texttt{eval} 的递归调用后，
这个调用尽管看起来像尾调用，却不再是尾调用。较好（但可读性稍差）的写法是：

\begin{minipage}{\linewidth}
\begin{taplocamlcode}
let rec eval t =
  let t'opt =
    try Some (eval1 t) with NoRuleApplies -> None
  in
  match t'opt with
    Some t' -> eval t'
  | None -> t
\end{taplocamlcode}
\taplrustcounterpart{sol-author-04-eval}
\end{minipage}

\taplsolutionbacklink{ex:4.2.1}{返回练习 4.2.1}
\end{taplauthorsolutionentrybox}
